\documentclass[12pt,letterpaper]{article}
\usepackage[utf8]{inputenc}
\usepackage[T1]{fontenc}
\usepackage[margin=2.54cm]{geometry}
\usepackage{mathptmx}
\usepackage[table]{xcolor}
\usepackage{array}
\usepackage[hidelinks]{hyperref}
\pagestyle{empty}
\color{black}
\renewcommand{\arraystretch}{1.25}
\definecolor{cab}{gray}{0.75}
\newcolumntype{A}{>{\raggedright\arraybackslash}p{3.1cm}}
\newcolumntype{B}{>{\raggedright\arraybackslash}p{7.4cm}}
\newcolumntype{C}{>{\raggedright\arraybackslash}p{4.2cm}}
\begin{document}
\begin{titlepage}
\centering
\vspace*{1cm}
{\Large \textbf{UNIVERSIDAD DE CÓRDOBA}}\\[0.3cm]
{\large FACULTAD DE INGENIERÍA}\\[0.2cm]
{\large PROGRAMA DE INGENIERÍA DE SISTEMAS}\\[0.2cm]
{\large SEDE MONTERÍA}\\[1.7cm]
\rule{\textwidth}{0.5pt}\\[0.3cm]
{\Huge \textbf{FICHAS DE HARDWARE Y SOFTWARE}}\\[0.4cm]
{\LARGE \textbf{SolarBell}}\\[0.3cm]
\rule{\textwidth}{0.5pt}\\[0.5cm]
{\large \textit{Fichas técnicas de los componentes y las funcionalidades del sistema SolarBell.}}\\[2.5cm]
\begin{minipage}[t]{0.42\textwidth}
\begin{flushleft}
\textbf{Presentado por:}\\[0.2cm]
Andrés Felipe Ortiz Rodríguez,\\
Néstor Junior Hernández López,\\
José Enrique Otero Pineda,\\
Julián Camilo Doria Guevara
\end{flushleft}
\end{minipage}\hfill
\begin{minipage}[t]{0.52\textwidth}
\begin{flushright}
\textbf{Docentes:}\\[0.2cm]
Antonio Fernando Movilla Quintero\\
\textit{Sistemas Operativos}\\[0.3cm]
Milton Eduardo Hernández Zakzuk\\
\textit{Arquitectura del Computador}
\end{flushright}
\end{minipage}\\[2.5cm]
\vfill
Montería, Córdoba\\
Octubre de 2026
\end{titlepage}
\newpage
\noindent{\large \textbf{Ficha de Hardware}}\\[0.25cm]
\noindent\begin{tabular}{|A|B|C|}
\hline
\rowcolor{cab}\multicolumn{2}{|l|}{\textbf{Nombre del Proyecto:}} & \textbf{Plataforma:} \\
\hline
\multicolumn{2}{|l|}{SolarBell: Sistema de timbre automatizado con energía solar} & ESP32 (con conectividad Wi-Fi) \\
\hline
\rowcolor{cab}\textbf{Funcionalidad} & \textbf{Descripción} & \textbf{Dispositivo} \\
\hline
Conectividad & El dispositivo (ESP32) crea su propia red Wi-Fi y, de forma opcional, se une a la red Wi-Fi local para recibir los horarios y las órdenes de la aplicación. & ESP32 DevKit v1 (Wi-Fi integrado en la plataforma ESP32) \\
\hline
Control de Tiempo & Mantiene la fecha y la hora exactas, incluso sin internet o sin energía, para activar el timbre en los horarios programados y en las fechas especiales. & Módulo RTC DS3231 (HW-084) \\
\hline
Activación del Inversor & Conecta la alimentación de 12 V del inversor solo cuando va a sonar la campana, para evitar su consumo en vacío y cuidar la batería. & Módulo de relé de 1 canal, 5 V (azul) \\
\hline
Activación de la Campana & Interrumpe y cierra la línea de 110 VAC que va hacia la campana para que suene en los horarios programados o en un toque manual. & Relé Songle SRA-05VDC-CL, 5 V \\
\hline
Activación Manual & Permite que la campana suene directamente al oprimir un botón, sin depender de la aplicación. & Pulsador de pared para timbre (EXE) \\
\hline
Generación de Energía & Convierte la radiación solar en energía eléctrica para alimentar el sistema. & Panel solar Powest SZYL-P30-18C (30 Wp) \\
\hline
Regulación de Carga & Regula la carga de la batería desde el panel y la protege contra sobrecarga y descarga profunda. & Controlador de carga solar PWM de 10 A (12/24 V) \\
\hline
Almacenamiento de Energía & Almacena la energía para que el sistema funcione de noche y en días sin sol. & Batería VRLA 12 V / 7.5 Ah (FuliBattery FL1275GS) \\
\hline
Conversión de Voltaje (DC-AC) & Convierte los 12 VDC de la batería en 110 VAC para alimentar la campana. & Inversor CAR INVERTER LY-150 (150 W) \\
\hline
Reducción de Voltaje & Convierte la energía de la batería en una salida regulada de 5 V destinada al ESP32 y a los módulos de relé. & Convertidor reductor LM2596 \\
\hline
\end{tabular}

\newpage

\noindent{\large \textbf{Ficha de Software}}\\[0.25cm]
\noindent\begin{tabular}{|A|B|C|}
\hline
\rowcolor{cab}\multicolumn{2}{|l|}{\textbf{Nombre del Proyecto:}} & \textbf{Plataforma:} \\
\hline
\multicolumn{2}{|l|}{SolarBell: Sistema de timbre automatizado con energía solar} & Móvil (solo Android) \\
\hline
\rowcolor{cab}\textbf{Funcionalidad} & \multicolumn{2}{l|}{\textbf{Descripción}} \\
\hline
Conectividad & \multicolumn{2}{>{\raggedright\arraybackslash}p{11.3cm}|}{El software se comunica con el dispositivo de hardware (ESP32) a través de su API por Wi-Fi, ya sea en la red propia del dispositivo (SolarBell) o en la red local, para consultar su estado y enviarle órdenes.} \\
\hline
Gestión de Horarios & \multicolumn{2}{>{\raggedright\arraybackslash}p{11.3cm}|}{Permite consultar y modificar el horario de timbres, con la duración de los toques y las pausas. El horario se guarda en la memoria del dispositivo para que funcione sin conexión.} \\
\hline
Fechas Especiales & \multicolumn{2}{>{\raggedright\arraybackslash}p{11.3cm}|}{Permite configurar fechas especiales (como festivos o días sin clases) en las que el horario normal de timbres se modifica o no se aplica.} \\
\hline
Toque Manual & \multicolumn{2}{>{\raggedright\arraybackslash}p{11.3cm}|}{Permite hacer sonar la campana en el momento que el usuario lo requiera y detener el toque en curso.} \\
\hline
Estado del Dispositivo & \multicolumn{2}{>{\raggedright\arraybackslash}p{11.3cm}|}{Permite consultar desde la aplicación el estado actual del dispositivo.} \\
\hline
\end{tabular}
\end{document}
